\documentclass[a4,12pt]{article}
\usepackage[utf8]{inputenc}
\usepackage[T2A]{fontenc}
\usepackage[english, russian]{babel}

\sloppy % Все тексты в документе будут верстаться с менее строгими правилами

\usepackage{csquotes}
\usepackage{epigraph}
\usepackage{tcolorbox}
% Определяем окружение для Формата ввода
\newtcolorbox{inputformat}[1][]{colback=blue!5!white, colframe=blue!65!black,
    fonttitle=\bfseries, title=Формат ввода, #1}

% Определяем окружение для Формата вывода
\newtcolorbox{outputformat}[1][]{colback=green!5!white, colframe=green!65!black,
    fonttitle=\bfseries, title=Формат вывода, #1}

% Определяем окружение для комментария
\newtcolorbox{exercisecomment}[1][]{colback=yellow!5!white, colframe=yellow!80!black, fonttitle=\bfseries, title=Комментарий, #1}

% Определяем окружение для замечания
%\newtcolorbox{exercisenote}[1][]{colback=red!5!white, colframe=red!80!black, fonttitle=\bfseries, title=Замечание, #1}

\newtcolorbox{exercisenote}[1][]{colback=red!5!white, colframe=red!80!black, fonttitle=\bfseries, title=Замечание, #1}
\newtcolorbox{exercisehint}[1][]{colback=gray!5!white, colframe=gray!60!black, fonttitle=\bfseries, title=Подсказка, #1}
\newtcolorbox{exerciseexample}[1][]{colback=blue!5!white, colframe=blue!65!black, fonttitle=\bfseries, title=Пример, #1}



\usepackage{geometry}
%\geometry{papersize={20.3 cm,25.4 cm}}
\geometry{left=2cm}
\geometry{right=2cm}
\geometry{top=2cm}
\geometry{bottom=2cm}


\usepackage{amsmath, amsfonts, amsthm, amssymb, amsopn, amscd}
\usepackage{enumerate}
\usepackage{enumitem}
\usepackage[mathscr]{eucal}

\usepackage{hyperref}
\hypersetup{unicode=true,final=true,colorlinks=true}

\theoremstyle{remark}
%\newtheorem{exercises}{Упражнения}
\newtheorem{exercise}{\textbf{Упражнение}}[section]
\renewcommand{\theexercise}{\textbf{\arabic{exercise}}}
%\renewcommand{\theexercise}{\textbf{\#\arabic{exercise}}}


\title{Листок 09. Включения в Python}

\author{А.Н. Лата}

\begin{document} 

%\maketitle

\section*{\centering Листок 09. Включения в Python}

\begin{exercisenote}[title=Замечания]
\begin{itemize}
    \item для решения приведённых ниже упражнений не требуется создавать (определять) пользовательские функции и использовать сторонние библиотеки.
    \item при решении упражнений полученные результаты выведите на экран с помощью функции {\color{blue}{print()}}.
    \item позаботесь о том, чтобы выводимый на экран результат был снабжен информацией о нем, там где это необходимо.
    \item не забывайте писать комментарии к вашему коду
\end{itemize}
\end{exercisenote}

\section*{\centering Упражнения} %Задачи на списковые, множественные и словарные включения

% =========================================================
\section*{Часть I. Списковые включения (list comprehensions)}
% =========================================================

\begin{exercisenote}[title=Синтаксис]
Списковое включение — это компактный способ создать список на основе
существующей последовательности. Общая форма записи:
\begin{center}
    \texttt{[выражение\ for\ элемент\ in\ последовательность\ if\ условие]}
\end{center}
Часть \texttt{if~условие} является необязательной и используется для фильтрации
элементов. Результатом всегда является список (\texttt{list}), порядок элементов
соответствует порядку обхода последовательности.
\end{exercisenote}

\begin{enumerate}
    \subsection*{Начальный уровень}
    \item Создайте список квадратов чисел от 1 до 10.
    \item Создайте список только чётных чисел от 0 до 20.
    \item Дан список строк \texttt{['hello', 'python', 'code']}. Создайте новый список, где каждая строка продублирована (например, из \texttt{'hello'} получается \texttt{'hellohello'}).
    \item Дан список слов \texttt{['python', 'java', 'javascript', 'css']}. Создайте список, содержащий длину каждого слова.
    \item Дано слово \texttt{'python'}. Создайте список, состоящий из букв этого слова в верхнем регистре.

    \subsection*{Средний уровень}
    \item Создайте список кортежей, где каждый кортеж содержит число и его квадрат для чисел от 1 до 5.
    \item Дан список строк \texttt{['cat', 'dog', 'elephant', 'bee', 'butterfly']}. Создайте новый список, содержащий только те строки, длина которых больше 3 символов.
    \item Создайте список всех чисел от 1 до 30, которые делятся либо на 3, либо на 5.
    \item Дан список температур по Цельсию \texttt{[0, 10, 20, 25, 30]}. Преобразуйте их в температуры по Фаренгейту.
    \item Дано предложение \texttt{'Python is a great programming language'}. Создайте список из первых букв каждого слова (только для непустых слов).

    \subsection*{Продвинутый уровень}
    \item Дана строка \texttt{'python'}. Создайте список всех возможных подстрок длины 2.

    \begin{exercisehint}
        Подстрока длины 2, начинающаяся с позиции \texttt{i}, записывается как \texttt{s[i:i+2]}.
        Перебирайте допустимые значения \texttt{i} в спецификации включения.
    \end{exercisehint}

    \item Дан список дат $\texttt{['2024-03-01', '2024-03-15', '2024-04-01']}$ в формате \texttt{'YYYY-MM-DD'}. Преобразуйте их в список кортежей \texttt{(год, месяц, день)}.
    \item Дано предложение \texttt{'Python programming is very interesting'}. Создайте список слов, которые содержат хотя бы одну гласную букву.
    \item Найдите все числа-палиндромы в диапазоне от 100 до 1000.
    \item Дан список чисел \texttt{[1, 2, 3, 4, 5]}. Создайте список, содержащий суммы каждой пары соседних элементов.
    \item Дан список строк \texttt{['contact@example.com', 'python.org', 'user@mail.com', 'no-reply@site.com', 'invalid.email']}.  
    Извлеките из него все адреса электронной почты (строки, содержащие символ \texttt{'@'}).

    \item Создайте список всех дат марта 2024 года в формате \texttt{'YYYY-MM-DD'}.
    (В марте 31 день; используйте \texttt{range(1, 32)} для перебора дней.)

    \item Три натуральных числа $a$, $b$ и $c$ называются \emph{пифагоровой тройкой},
    если они могут быть сторонами прямоугольного треугольника,
    то есть выполняется равенство $a^2 + b^2 = c^2$,
    где $c$ --- длина гипотенузы, $a$ и $b$ --- длины катетов.
    Например, тройка $(3, 4, 5)$: $3^2 + 4^2 = 9 + 16 = 25 = 5^2$.

    % Найдите все пифагоровы тройки $(a,\, b,\, c)$, в которых каждое из чисел
    % не превышает 20. Чтобы каждая тройка встречалась в ответе ровно один раз,
    % соблюдайте условие $a < b < c$.

    Найдите все пифагоровы тройки $(a,\, b,\, c)$ при условии $a < b < c \leqslant 20$.

\end{enumerate}

% =========================================================
\section*{Часть II. Множественные включения (set comprehensions)}
% =========================================================

\begin{exercisenote}[title=Синтаксис]
Множественное включение записывается аналогично списковому включению,
но в фигурных скобках:
\begin{center}
    \texttt{\{выражение\ for\ элемент\ in\ последовательность\ if\ условие\}}
\end{center}
Результатом является множество (\texttt{set}), поэтому порядок элементов
не гарантирован, а дубликаты автоматически устраняются.
\end{exercisenote}


\begin{enumerate}

    \subsection*{Начальный уровень}

    \item Создайте множество квадратов чисел от 1 до 10.

    \begin{exerciseexample}
        \texttt{\{1, 4, 9, 16, 25, 36, 49, 64, 81, 100\}}
    \end{exerciseexample}

    \item Дан список \texttt{['apple', 'banana', 'apple', 'orange', 'banana', 'apple']}.
    Создайте множество уникальных элементов этого списка с помощью
    множественного включения.

    \item Создайте множество всех чётных чисел от 0 до 20.

    \item Дана строка \texttt{'программирование'}.
    Создайте множество уникальных букв этой строки.

    \item Дан список чисел \texttt{[1, 2, 3, 4, 5, 6, 7, 8, 9]}.
    Создайте множество остатков от деления каждого числа на 3.

    \subsection*{Средний уровень}

    \item Дан список слов \texttt{['python', 'java', 'javascript', 'python', 'css', 'kotlin']}.
    Создайте множество слов длиной более 4 символов
    (дубликаты должны присутствовать только один раз).

    \item Создайте множество всех нечётных чисел от 1 до 50, кратных 3.

    \item Дана строка \texttt{'Hello, World!'}.
    Создайте множество уникальных строчных букв этой строки,
    исключив пробелы и знаки препинания.

    \begin{exercisehint}
        Метод \texttt{str.isalpha()} возвращает \texttt{True},
        если символ является буквой.
    \end{exercisehint}

    \item Даны два списка \texttt{a = [1, 2, 3, 4, 5]} и \texttt{b = [3, 4, 5, 6, 7]}.
    Создайте с помощью одного множественного включения множество всех элементов,
    входящих хотя бы в один из этих двух списков (объединение).

    \begin{exercisehint}
        Объедините оба списка операцией \texttt{a + b} перед или внутри включения.
    \end{exercisehint}

    \item Дано предложение \texttt{'Python is a great programming language'}.
    Создайте множество первых букв слов этого предложения.

    \subsection*{Продвинутый уровень}

    \item Создайте множество всех простых чисел до 50 с помощью
    множественного включения.

    \begin{exercisehint}
        Число \texttt{n} является простым, если ни одно число от 2 до
        \texttt{n - 1} не делит его нацело. Условие простоты можно записать
        с помощью \texttt{all()} внутри включения.
    \end{exercisehint}

    \item Дан список строк \texttt{['cat', 'dog', 'cat', 'bird', 'dog', 'fish', 'dog']}.
    Создайте множество строк, которые встречаются в списке \emph{более одного раза}.

    \begin{exercisehint}
        Метод \texttt{list.count(x)} возвращает количество вхождений \texttt{x}
        в список.
    \end{exercisehint}

    \item Дан список чисел \texttt{[12, 45, 7, 33, 20, 18, 9, 100, 66]}.
    Создайте множество, содержащее суммы цифр каждого числа из списка.

    \begin{exercisehint}
        Цифры числа \texttt{n} можно получить, перебирая символы строки \texttt{str(n)}
        и преобразуя каждый символ в \texttt{int}.
    \end{exercisehint}

    \item Дана строка \texttt{'aabbccaabddee'}.
    Создайте множество символов, которые встречаются в строке чётное число раз.

\end{enumerate}


% =========================================================
\section*{Часть III. Словарные включения (dict comprehensions)}
% =========================================================

\begin{exercisenote}[title=Синтаксис]
Словарное включение записывается в фигурных скобках с явным указанием
пары ключ\,:\,значение:
\begin{center}
    \texttt{\{ключ: значение\ for\ элемент\ in\ последовательность\ if\ условие\}}
\end{center}
Результатом является словарь (\texttt{dict}).
\end{exercisenote}

\begin{enumerate}

    \subsection*{Начальный уровень}

    \item Создайте словарь, где ключи — числа от 1 до 5, а значения — их квадраты.

    \begin{exerciseexample}
        \texttt{\{1: 1,\ 2: 4,\ 3: 9,\ 4: 16,\ 5: 25\}}
    \end{exerciseexample}

    \item Дан список слов \texttt{['python', 'java', 'css']}.
    Создайте словарь, где ключи — слова, а значения — длины этих слов.

    \item Создайте словарь, где ключи — буквы строки \texttt{'abcde'},
    а значения — их порядковые номера, начиная с 1.

    \begin{exercisehint}
        Используйте \texttt{enumerate()} в спецификации включения.
    \end{exercisehint}

    \item Дан словарь \texttt{\{'a': 1, 'b': 2, 'c': 3\}}.
    Создайте новый словарь, в котором ключи и значения поменяны местами.

    \item Дан список чисел \texttt{[1, 2, 3, 4, 5]}.
    Создайте словарь, где ключи — числа, а значения — их кубы.

    \subsection*{Средний уровень}

    \item Дан список слов \texttt{['apple', 'banana', 'cherry', 'fig', 'date']}.
    Создайте словарь, содержащий только слова длиной более 4 символов
    (ключ — слово, значение — его длина).

    \item Создайте словарь, где ключи — числа от 1 до 10,
    а значения — строка \texttt{'чётное'} или \texttt{'нечётное'}.

    \item Дан список пар \texttt{[('a', 1), ('b', 2), ('c', 3), ('d', 4)]}.
    Создайте словарь из этих пар с помощью словарного включения.

    \item Даны два списка \texttt{keys = ['x', 'y', 'z']} и
    \texttt{values = [10, 20, 30]}.
    Создайте словарь, объединив их попарно.

    \begin{exercisehint}
        Функция \texttt{zip(keys, values)} создаёт последовательность
        пар \texttt{(ключ, значение)}, по которой удобно итерировать.
    \end{exercisehint}

    \item Дана строка \texttt{'python'}.
    Создайте словарь, где ключи — буквы строки,
    а значения — индексы их первого появления.

    \begin{exercisehint}
        Метод \texttt{str.index(символ)} возвращает индекс первого вхождения
        символа в строку.
    \end{exercisehint}

    \item Дан словарь \texttt{\{'a': 3, 'b': 1, 'c': 4, 'd': 2, 'e': 5\}}.
    Создайте новый словарь, содержащий только те пары,
    у которых значение больше 2.

    \item Создайте словарь, где ключи — числа от 1 до 15,
    а значения — строка \texttt{'FizzBuzz'} (если делится на 3 и на 5),
    \texttt{'Fizz'} (если делится только на 3), \texttt{'Buzz'}
    (если делится только на 5) или само число (в остальных случаях).

    \begin{exercisehint}
        Используйте вложенные тернарные выражения:
        \texttt{A if условие1 else B if условие2 else C}.
    \end{exercisehint}

    \subsection*{Продвинутый уровень}

    \item Дан текст \texttt{'python is great and python is easy'}.
    Создайте словарь частот слов: ключ — слово, значение — количество его вхождений.

    \begin{exercisehint}
        Метод \texttt{str.split()} разбивает строку на список слов.
        Метод \texttt{list.count(x)} возвращает количество вхождений элемента.
    \end{exercisehint}

    \item Дан список строк \texttt{['apple', 'banana', 'avocado', 'blueberry', 'cherry', 'apricot']}.
    Создайте словарь, где ключи — первые буквы слов,
    а значения — списки слов, начинающихся на эту букву.

    \begin{exercisehint}
        Используйте список в качестве значения: для каждой уникальной первой буквы
        \texttt{c} отберите с помощью спискового включения все слова,
        у которых \texttt{word[0] == c}.
    \end{exercisehint}

    \item Создайте словарь, где ключи — числа от 1 до 10,
    а значения — списки всех их делителей.

    \begin{exerciseexample}
        \texttt{\{1: [1],\ 2: [1, 2],\ 3: [1, 3],\ 4: [1, 2, 4],\ \ldots\}}
    \end{exerciseexample}

    \item Дан список чисел \texttt{[1, 2, 3, 4, 5, 6, 7, 8, 9, 10]}.
    Создайте словарь, где ключи — элементы списка,
    а значения — их квадраты; при этом включайте в словарь
    только те числа, квадрат которых оканчивается на 1, 4, 5, 6 или 9.

    \item Дан текст \texttt{'Hello World from Python'}.
    Создайте словарь, где ключи — слова текста,
    а значения — словари вида \texttt{\{'длина': ..., 'заглавная': ...\}},
    содержащие длину слова и его вариант с заглавной буквы.

    \begin{exerciseexample}
        \texttt{\{'Hello': \{'длина': 5, 'заглавная': 'Hello'\},\ \ldots\}}
    \end{exerciseexample}

\end{enumerate}

\end{document}